\documentclass{book}
\makeatletter
\def\partbody{This partition paragraph exercises the typesetting pipeline with mixed content including mathematical expressions such as $\alpha_{n} + \beta_{n} = \gamma_{n}$ and the recurrence $a_{n+1} = a_n + n$ to establish meaningful compilation workload for parallel benchmarking. A second sentence revisits deterministic page construction, numbering, and paragraph breaking while mentioning the summation $\sum_{k=1}^{n} k = \frac{n(n+1)}{2}$ and the ratio $\frac{s_n}{n + 1} = n$. A third sentence keeps the content stable across job counts by repeating predictable prose about counters, references, and layout decisions so each partition contributes measurable compilation work. }
\def\partworkload{%
  \@for\pitem:=01,02,03,04,05,06,07,08,09,10,11,12,13,14,15,16,17,18,19,20,21,22,23,24,25,26,27,28,29,30,31,32,33,34,35,36,37,38,39,40,41,42,43,44,45,46,47,48,49,50\do{%
    \partbody
    \partbody
    \[
    \frac{s_n}{n + 1} = n
    \]
    \begin{equation}
    s_n = n^2 + n
    \end{equation}
  }%
  \@for\pitem:=051,052,053,054,055,056,057,058,059,060,061,062,063,064,065,066,067,068,069,070,071,072,073,074,075,076,077,078,079,080,081,082,083,084,085,086,087,088,089,090,091,092,093,094,095,096,097,098,099,100\do{%
    \partbody
    \partbody
    \[
    \frac{s_n}{n + 1} = n
    \]
    \begin{equation}
    s_n = n^2 + n
    \end{equation}
  }%
}
\makeatother
\begin{document}
\chapter{Chapter One}\label{chap:one}
Chapter one introduces the references to Chapter \ref{chap:ten}.
\makeatletter
\partworkload
\makeatother
\newpage

\chapter{Chapter Two}\label{chap:two}
Chapter two points back to Chapter \ref{chap:one}.
\makeatletter
\partworkload
\makeatother
\newpage

\chapter{Chapter Three}\label{chap:three}
Chapter three references Chapter \ref{chap:two}.
\makeatletter
\partworkload
\makeatother
\newpage

\chapter{Chapter Four}\label{chap:four}
Chapter four references Chapter \ref{chap:three}.
\makeatletter
\partworkload
\makeatother
\newpage

\chapter{Chapter Five}\label{chap:five}
Chapter five references Chapter \ref{chap:four}.
\makeatletter
\partworkload
\makeatother
\newpage

\chapter{Chapter Six}\label{chap:six}
Chapter six references Chapter \ref{chap:five}.
\makeatletter
\partworkload
\makeatother
\newpage

\chapter{Chapter Seven}\label{chap:seven}
Chapter seven references Chapter \ref{chap:six}.
\makeatletter
\partworkload
\makeatother
\newpage

\chapter{Chapter Eight}\label{chap:eight}
Chapter eight references Chapter \ref{chap:seven}.
\makeatletter
\partworkload
\makeatother
\newpage

\chapter{Chapter Nine}\label{chap:nine}
Chapter nine references Chapter \ref{chap:eight}.
\makeatletter
\partworkload
\makeatother
\newpage

\chapter{Chapter Ten}\label{chap:ten}
Chapter ten closes the book and references Chapter \ref{chap:one}.
\makeatletter
\partworkload
\makeatother
\end{document}
